% BE-04: System Context Diagram — compile with XeLaTeX (Overleaf: Menu → Compiler → XeLaTeX)
\documentclass[11pt,a4paper]{article}

% ---------- language and fonts ----------
\usepackage{fontspec}
\usepackage{polyglossia}
\setdefaultlanguage{english}
% fonts with Cyrillic glyphs, needed for the Kazakh/Russian command examples
\IfFontExistsTF{Times New Roman}{\setmainfont{Times New Roman}}{%
  \IfFontExistsTF{Liberation Serif}{\setmainfont{Liberation Serif}}{\setmainfont{DejaVu Serif}}}
\IfFontExistsTF{Arial}{\setsansfont{Arial}}{%
  \IfFontExistsTF{Liberation Sans}{\setsansfont{Liberation Sans}}{\setsansfont{DejaVu Sans}}}
\IfFontExistsTF{Consolas}{\setmonofont[Scale=0.9]{Consolas}}{\setmonofont[Scale=0.85]{DejaVu Sans Mono}}

% ---------- layout ----------
\usepackage[margin=2.2cm]{geometry}
\usepackage{microtype}
\usepackage{array}
\usepackage{booktabs}
\usepackage{xltabular}
\usepackage{enumitem}
\usepackage[table]{xcolor}
\usepackage{caption}
\usepackage{float}
\usepackage{tikz}
\usetikzlibrary{positioning,arrows.meta,fit,backgrounds}
\usepackage{xurl}
\usepackage[colorlinks=true,linkcolor=blue!50!black,urlcolor=blue!50!black]{hyperref}

\setlist{itemsep=2pt,topsep=4pt}
\setlength{\parskip}{4pt}
\setlength{\parindent}{0pt}
\setlength{\emergencystretch}{3em}
\renewcommand{\arraystretch}{1.25}
\newcolumntype{L}{>{\raggedright\arraybackslash}X}
\newcolumntype{P}[1]{>{\raggedright\arraybackslash}p{#1}}
\newcommand{\tblstart}{\small\rowcolors{2}{gray!7}{white}}
\newcommand{\code}[1]{\texttt{#1}}

% colours of the C4 notation
\definecolor{c4person}{HTML}{08427B}
\definecolor{c4system}{HTML}{1168BD}
\definecolor{c4ext}{HTML}{8A8A8A}

\title{\textbf{BE-04: System Context Diagram}\\[4pt]
\large Intelligent Kanban System with an AI Assistant}
\author{Symbat Saparbay | Akbope Bakytkeldy\\ \small YouTrack: DEMO-23}
\date{}

\begin{document}
\maketitle
% =====================================================
\section{Purpose}

The context diagram shows the system as a single whole: \textbf{who} uses it, \textbf{which external systems} it interacts with, and \textbf{what data} crosses its boundary. The internal structure (frontend, backend, AI worker, databases) is described in BE-05 (Backend Component Diagram).


% =====================================================
\section{Context Diagram}

\begin{figure}[H]
\centering
\begin{tikzpicture}[
  font=\small,
  person/.style={draw=c4person, fill=c4person, text=white, rounded corners=8pt, align=center,
                 text width=3.2cm, minimum height=1.4cm, inner sep=4pt},
  system/.style={draw=c4system, fill=c4system, text=white, rounded corners=6pt, align=center,
                 text width=9.6cm, minimum height=2.6cm, inner sep=7pt},
  ext/.style={draw=c4ext, fill=c4ext, text=white, rounded corners=6pt, align=center,
              text width=3.6cm, minimum height=1.4cm, inner sep=4pt},
  extopt/.style={ext, dashed, draw=c4ext!80!black, fill=c4ext!25, text=black!85},
  arr/.style={-{Stealth[length=2.2mm]}, thick, draw=gray!60!black},
  tag/.style={font=\scriptsize, fill=white, inner sep=2pt, align=center, text=black}
]
% people (top)
\node[person] (owner)   at (-6.3, 4.6) {\textbf{Workspace}\\\textbf{owner / admin}\\{\footnotesize members, roles, settings}};
\node[person] (manager) at (-2.1, 4.6) {\textbf{Project}\\\textbf{manager}\\{\footnotesize manager, teacher}};
\node[person] (member)  at ( 2.1, 4.6) {\textbf{Project}\\\textbf{member}\\{\footnotesize employee, student}};
\node[person] (viewer)  at ( 6.3, 4.6) {\textbf{Viewer}\\[2pt]{\footnotesize read-only access}};

% system (centre)
\node[system] (sys) at (0,0) {%
  {\large\bfseries Intelligent Kanban System}\\[5pt]
  Kanban boards, tasks, comments, history, notifications.\\
  The AI assistant understands voice and text commands in\\
  Kazakh, Russian and English and proposes actions\\
  that are performed after confirmation.\\[4pt]
  {\footnotesize Speech recognition (Whisper) and the local LLM run inside the system}};

% external systems (bottom)
\node[ext]    (llm)  at (-5.2,-4.4) {\textbf{Cloud LLM API}\\{\footnotesize Gemini Flash / OpenAI mini}};
\node[ext]    (mail) at ( 0,  -4.4) {\textbf{Email service}\\{\footnotesize SMTP}};
\node[extopt] (stt)  at ( 5.2,-4.4) {\textbf{Cloud STT}\\{\footnotesize fallback, disabled}\\{\footnotesize by default}};

% people -> system
\draw[arr, <->] (owner.south)   -- node[tag]{manages members\\and AI settings} ([xshift=0.9cm]sys.north west);
\draw[arr, <->] (manager.south) -- node[tag]{projects, tasks,\\voice commands} (manager.south |- sys.north);
\draw[arr, <->] (member.south)  -- node[tag]{own tasks,\\statuses, comments} (member.south |- sys.north);
\draw[arr]      (viewer.south)  -- node[tag]{views\\boards} ([xshift=-0.9cm]sys.north east);

% system -> external
\draw[arr, <->] ([xshift=0.9cm]sys.south west) -- node[tag]{command text $\rightarrow$\\$\leftarrow$ action plan (JSON)} (llm.north);
\draw[arr] (sys.south) -- node[tag]{notifications,\\deadline reminders} (mail.north);
\draw[arr, <->, dashed] ([xshift=-0.9cm]sys.south east) -- node[tag]{audio $\rightarrow$\\$\leftarrow$ transcript} (stt.north);

% legend
\node[anchor=north, font=\scriptsize, align=center] at (0,-5.6) {%
  \tikz\node[fill=c4person, minimum width=0.5cm, minimum height=0.3cm, rounded corners=2pt]{}; Person\qquad
  \tikz\node[fill=c4system, minimum width=0.5cm, minimum height=0.3cm, rounded corners=2pt]{}; Our system\qquad
  \tikz\node[fill=c4ext, minimum width=0.5cm, minimum height=0.3cm, rounded corners=2pt]{}; External system\qquad
  \tikz\node[draw=c4ext!80!black, dashed, fill=c4ext!25, minimum width=0.5cm, minimum height=0.3cm, rounded corners=2pt]{}; Optional external system};
\end{tikzpicture}
\caption{System context diagram. Users work through a web browser over HTTPS/WSS; emails from the email service reach users' mailboxes.}
\label{fig:context}
\end{figure}

% =====================================================
\section{Users}

The roles follow the RBAC model from BE-03: workspace roles (owner, admin, member) and project roles (manager, member, viewer). One person can have different roles in different projects.

{\tblstart
\begin{xltabular}{\textwidth}{P{3.1cm} L L}
\toprule
\textbf{Role} & \textbf{What they do in the system} & \textbf{Example AI commands} \\
\midrule
\endhead
Workspace owner / admin & Creates the workspace, invites members, assigns roles; enables or disables the cloud LLM (private mode) and the cloud STT fallback; views the audit log & None (managed through the UI) \\
Project manager & Creates projects and boards, adds members, creates and assigns tasks, sets deadlines and priorities, records meeting outcomes & «Создай задачу для Данияра: подготовить презентацию до пятницы, высокий приоритет»; “Which of our team's tasks are overdue?” \\
Project member & Works on their own tasks: changes status, comments, creates tasks in accessible projects; receives notifications & «Задача по интеграции платёжной системы завершена»; «Даниярға ертеңге дейін API тексеру тапсырмасын бер» \\
Viewer & Only views boards and tasks (for example, a teacher or a client); cannot change data & Search Tasks and Project Summary only \\
\bottomrule
\end{xltabular}}

% =====================================================
\section{External Systems}

{\tblstart
\begin{xltabular}{\textwidth}{P{2.8cm} P{2.4cm} L P{3.2cm}}
\toprule
\textbf{System} & \textbf{Required?} & \textbf{Interaction} & \textbf{If unavailable} \\
\midrule
\endhead
Cloud LLM API (Gemini Flash / OpenAI mini) & Yes in the default mode; no in private mode & Receives the command text (after speech recognition), project member names, the current date and time zone. Returns an action plan in JSON: action type, parameters, missing fields & Switch to the local LLM (Ollama); manual task management always works \\
Email service (SMTP) & Yes & Sends emails: task assignment, mention, upcoming and overdue deadline, workspace invitation & Emails are queued and retried; in-app notifications keep working \\
Cloud STT (ElevenLabs Scribe / OpenAI) & No, disabled by default & Only if an admin enables the fallback: receives audio, returns a transcript & The local STT is used; if it is also unavailable, the user types the command \\
\bottomrule
\end{xltabular}}

\textbf{Why speech recognition is not an external system.} Whisper (including the Kazakh model) and the local LLM (Ollama) are deployed on our server and are part of the system. Therefore, in the default mode, voice \textbf{does not leave the system boundary}.

% =====================================================
\section{Data Crossing the System Boundary}

{\tblstart
\begin{xltabular}{\textwidth}{P{3.4cm} P{2.6cm} L}
\toprule
\textbf{Data} & \textbf{Sent to} & \textbf{Protection} \\
\midrule
\endhead
Voice messages & Nowhere (default) & Stored in MinIO with encryption, accessed via presigned URLs, deleted after 30 days. Sent to the cloud STT only if an admin enables the fallback \\
Command text & Cloud LLM API & HTTPS; only the text of a single command is sent, without audio or the contents of other tasks. Not sent anywhere in private mode \\
Project member names & Cloud LLM API & Only the names of members of the current project, without email addresses or other personal data; needed so the model can identify the assignee mention \\
User email and notification text & Email service & TLS; the email contains the task title and a link, without transcripts or audio \\
\bottomrule
\end{xltabular}}

Data that \textbf{never} leaves the system: passwords (stored as Argon2 hashes), tokens, the audit log, and the contents of tasks and comments from other projects.

% =====================================================
\section{Scope and Assumptions}

\begin{itemize}
  \item All users work through a web browser (desktop and mobile); voice is recorded in the browser. There is no separate mobile app.
  \item Sign-in is by email and password; external identity providers (Google, university SSO) are not part of the first version.
  \item Integrations with external task trackers, messengers and calendars (Jira, Telegram, Google Calendar) are not part of the first version; the architecture allows adding them later as new external systems.
  \item Development tools (GitHub/GitLab, YouTrack) are not part of the running system and are not shown on the diagram.
  \item The AI assistant never changes data on its own: every action is performed only after the user confirms it and only within the user's permissions.
\end{itemize}
\end{document}
